% NOETHER PAPER 2 — KOREAN TRANSLATION-PRODUCER DRAFT — UNCHECKED
% Unit: T08-U40; current-authority whole lines 1203--1210
% Source slice: 556 LF UTF-8 bytes; SHA-256 32BF9A8FCFCCA8001C4209776C5E6172F6F3D85B9EF1D8F713D617AE94BF7374
% Pointer: NOETH-DE-AUTH-v004-20260804; SHA-256 A1C62FDACAA34DFC1B806DC18258F2E732539F7AE9D85AA4BD9E1067B8749D9F
% This is producer translation only: no source/Korean/formula review, compilation, or rendering.

\bigskip

따라서 $\nu$를 법으로 하는 상대 완전계는 다음 여섯 형식으로 이루어진다.
\[
f,\ \theta,\ K,\ j,\ \A,\ N.
\]

§ 2의 급수 전개에 대한 예로 공식 (3.)a의 유도를 자세히 제시한다. 이후의 계산에서는 최종 공식 III을 직접 세울 것이다. (소멸하는 형식의 편극은 계산 도중에 이미 생략하였다. 첫째 줄은 전개 II에 따라 대입한 값을, 둘째 줄은 연립계의 마지막 방정식부터 시작하여 재귀적으로 구한 최종값을 나타낸다.)
